\documentclass[11pt,a4paper,sans]{moderncv}       
\moderncvstyle{classic}                          
\moderncvcolor{blue}                              
\usepackage[dvipsnames,svgnames]{xcolor}

\usepackage[normalem]{ulem}
\useunder{\uline}{\ul}{}
\usepackage[scale=0.78]{geometry}
\usepackage{CJKutf8}
\usepackage{fontawesome5}
\usepackage{simpleicons}
\usepackage{lmodern}
\usepackage[scaled]{helvet}

% --- Informations Personnelles ---
\name{Eiichiro}{Oda}
\title{\Large{Mangaka - Créateur de "One Piece"}}  

\born{1er Janvier 1975 (Kumamoto, Japon)}
\phone[mobile]{+81~03~5555~0000} 
\email{oda.eiichiro@shonenjump.co.jp}
\extrainfo{\faGlobe \href{https://one-piece.com/}{\;one-piece.com} \; {\scriptsize\simpleicon{shonenjump}} \href{https://www.shonenjump.com}{Shonen Jump}}

% Assure-toi d'avoir un fichier image nommé avatar.jpg dans le dossier
\photo[55pt][0.0pt]{avatar.jpg}

\begin{document}

\makecvtitle

\section{Expérience Professionnelle}

\cventry{Depuis 1997}{Auteur et Illustrateur}{Weekly Shōnen Jump}{Tokyo}{Japon}{
\begin{itemize}
    \item Création et sérialisation de la série \textbf{"One Piece"}, devenue le manga le plus vendu au monde.
    \item Gestion du world-building complexe, du design de plus de 1000 personnages et de l'intrigue sur plus de 1100 chapitres.
    \item Collaboration avec la Toei Animation pour l'adaptation animée et avec Netflix pour la série Live Action.
    \item Supervision des produits dérivés, films (Red, Stampede) et jeux vidéo.
\end{itemize}
}

\cventry{1994--1996}{Assistant Mangaka}{Shueisha}{}{}{
\begin{itemize}
    \item Assistant de \textbf{Nobuhiro Watsuki} sur le manga \textit{Rurouni Kenshin}.
    \item Assistant de Shinobu Kaitani (\textit{Midoriyama Police Gang}) et Masaya Tokuhiro (\textit{Jungle no Ouja Tar-chan}).
    \item Apprentissage des techniques d'encrage, de trames et de mise en page professionnelle.
\end{itemize}
}

\section{Éducation}
\cventry{1993}{Études en Architecture (Abandon)}{Université de Kyushu Sangyo}{Fukuoka}{Japon}{
\begin{itemize}
    \item Formation interrompue pour se consacrer pleinement à sa carrière de mangaka.
\end{itemize}
}
\cventry{1990--1993}{Lycée de Kumamoto}{Kumamoto}{Japon}{Spécialisation en dessin et arts plastiques.}{}

\section{Récompenses et Records}

\cvlistitem{\textbf{Guinness World Record (2015/2022) :} Série de bandes dessinées ayant le plus grand nombre d'exemplaires publiés par un seul auteur.}
\cvlistitem{\textbf{Grand Prix de l'Association des auteurs de bande dessinée japonais (2012) :} Pour l'excellence de l'œuvre One Piece.}
\cvlistitem{\textbf{Tezuka Award (1992) :} Finaliste à 17 ans pour son court récit \textit{Wanted!}.}
\cvlistitem{\textbf{Sondage Oricon :} Régulièrement élu "Mangaka le plus apprécié au Japon".}

\section{Compétences Techniques}
\cvitem{Dessin}{Encraque traditionnel (G-Pen), colorisation aux marqueurs Copic, storyboarding (Nems).}
\cvitem{Écriture}{Narration épique, construction de systèmes de combat (Fruits du Démon, Haki), gestion de la continuité narrative long-terme.}
\cvitem{Productivité}{Capacité de travail intense (cycle de sommeil de 3h par jour rapporté durant les périodes de production).}

\section{Œuvres Notables (One-shots)}
\cvlistitem{\textbf{Wanted!} (1992) - Son premier succès reconnu par le prix Tezuka.}
\cvlistitem{\textbf{Monsters} (1994) - Histoire centrée sur le samouraï Ryuma (lié plus tard à One Piece).}
\cvlistitem{\textbf{Romance Dawn} (1996) - Le prototype qui a servi de base à One Piece.}

\section{Langues}
\cvdoubleitem{Japonais}{Langue maternelle}{Anglais}{Notions de base}{}

\section{Centres d'intérêt}
\cvitem{Musique}{Fan de Soul music, de Brook Benton et d'Eminem (inspiration pour le design d'Enel).}
\cvitem{Divers}{Grand fan de l'histoire des pirates (Henry Morgan, Edward Teach) et de culture occidentale.}

\end{document}